%% 03 model details 
\subsection{Assume-Guarantee join operator + device Model}


\begin{enumerate}

    \item device components 
    \item operator component, an abstraction of human perception 
    \item assume and guarantees 
    \item sources from the cognitive literature for estimating the parameter values int he assumption, guarantee formulas.. 
    \item usage in designing filter, and perform security analysis.. 
    
\end{enumerate}

\subsection{HSR Protocol for Validting the Operator Component}

We plan conduct human subject experimentations to validate the operator component.  
To validate the operator model, we have devised the following Subject Research Protocol

Description: In an environment where the ambient lighting is constant, there is 1 target in the field of view.
Target walks along path starting from step 1 to 5. A bounding box around the target appears between step 2-
4. The timing of when a bounding box is randomized to prevent participant learning effect. The location of
the bounding box relative to the target is constant, but its appearance is also randomized to minimize
learning effect. The design choices of the bounding box incorporates current Helmet-Mounted symbology
concepts for degraded environments (Cheung et al., 2015).
Participants Instructions: Bounding boxes may or may not appear around a pedestrian. If you see a
bounding box, press a button as soon as you see it.
Measure: Participant performance in 

\begin{enumerate}
    \item true and false positives for identifying existence of bounding box
    \item delays in correct identification.
\end{enumerate}
Baseline Condition: No attack, target walks along path. Bounding box appears.
Attack Condition: Glare attack: Target turns on a headlamp somewhere along the walking path.
Luminance level of headlamp matches or is greater than bounding box. Headlamp is pointing at HMD.
Bounding box appears
Mitigation Condition: HMD modifies the stimuli by introducing a regional diming filter around the high
luminance area produced in the headlamp. See example of auto-diming glass
(https://www.gentextech.com/dimmable-glass.html).
Timing: Six combinations with an estimated 15 seconds per run. Each subject performs multiple sets of the
full combination for a total of 10 minute of subject task time. We may choose to add variations in ambient
lighting conditions should power analysis permits.
Hypotheses for Baseline, Attack, \& Mitigation conditions:
\begin{enumerate}
\item Baseline: Participants will correctly identify symbology (DV1) at $90\%$ with delays (DV2) <=600ms (+/- 60ms).
\item  Baseline DV1 estimates consider the presence and absence of factors specific to our
stimuli that affect human error rates (Remington \& Willaims, 1985).
\item Baseline DV2 estimates delay leverages findings in the analogous visual-motor task of a
cricket batter perceiving the ball release to judging the futural arrival point of the ball
(Sarpeshkar \& Mann, 2011).
\item Attack: Participants will correctly identify symbology (DV1) at $80\%$ with delays (DV2) >= 1sec.
\end{enumerate}
The expected performance decrement in delay is based on a hypothesized function of the following features:
\begin{enumerate} 
    \item Attack light source of 100 lumens such as from a low cost flashlight.
\item A maximum pupil size change of 1.15mm for a maximum luminance change of 10.90cd/m2 (standard deviation 20.04) (see Table 1 and 2, Sulutvedt et al., 2021)
\item 1mm increase in pupil diameter corresponds to an average increase in the perceptual experience of 2.09 cd/m2 in brightness (Sulutvedt et al., 2021)
\end{enumerate}
Mitigation: Return to baseline for both DV1 and DV2 with a $10\%$ performance decrement. 

Independent Variables

%% 